\documentclass{beamer}
\usepackage{graphicx}
\usepackage{booktabs}

% Theme and title setup
\usetheme{Frankfurt}
\title{Introduction to Contemporary Economic Issues in Canada and Globally}
\author{Muhebullah Karimzada}
\date{Winter 2025}

\begin{document}

% Title slide
\begin{frame}
  \titlepage
\end{frame}

% Slide: Course Overview
\begin{frame}{Course Overview}
  \begin{itemize}
    \item This course examines key economic challenges and opportunities in Canada and the global economy.
    \item Topics include globalization, inequality, technological change, climate economics, and post-pandemic recovery.
    \item The course integrates theory, empirical analysis, and policy discussions.
  \end{itemize}
\end{frame}

% Instructor Notes: Course Overview
\note{
Highlight that this course takes a practical approach to contemporary issues, integrating Canadian and global perspectives. Encourage students to think critically about the interconnectedness of these topics.
}

% Slide: Importance of Contemporary Economics
\begin{frame}{Why Study Contemporary Economics?}
  \begin{itemize}
    \item Understand the forces shaping modern economies.
    \item Analyze economic policies affecting society, businesses, and individuals.
    \item Address pressing global and national challenges, such as climate change and inequality.
    \item Develop critical thinking and data analysis skills applicable in policymaking and industry.
  \end{itemize}
\end{frame}

% Instructor Notes: Why Study Contemporary Economics
\note{
Explain that economics provides a lens to evaluate real-world problems and policy decisions. Use examples like recent climate policies or inflation trends to connect abstract concepts to real issues. Encourage students to bring in current news stories during discussions.
}

% Slide: Key Themes
\begin{frame}{Key Themes of the Course}
  \begin{itemize}
    \item \textbf{Globalization and Trade:} Impact on Canada and global economies.
    \item \textbf{Income Inequality:} Trends, causes, and policy responses.
    \item \textbf{Climate Change Economics:} Balancing growth with sustainability.
    \item \textbf{Technological Disruption:} Automation, AI, and labor markets.
    \item \textbf{Post-Pandemic Economic Recovery:} Fiscal and monetary challenges.
  \end{itemize}
\end{frame}

% Instructor Notes: Key Themes
\note{
Discuss how each theme is interconnected. For example, globalization impacts inequality, and technological disruption shapes recovery. Ask students to consider examples they’ve seen in the news or their own experiences.
}

% Slide: Globalization and Trade
\begin{frame}{Globalization and Trade}
  \begin{itemize}
    \item \textbf{Definition:} The increasing integration of economies through trade, investment, and technology.
    \item \textbf{Impacts on Canada:}
      \begin{itemize}
        \item Opportunities for export-driven growth (e.g., natural resources, technology).
        \item Challenges: Competition from low-cost producers and trade disputes.
      \end{itemize}
    \item \textbf{Global Perspective:}
      \begin{itemize}
        \item Rising protectionism (e.g., tariffs, trade wars).
        \item Role of international institutions like WTO, IMF.
      \end{itemize}
  \end{itemize}
\end{frame}

% Instructor Notes: Globalization and Trade
\note{
Use real-world examples, such as the impact of NAFTA/USMCA on Canada or recent trade disputes with China. Discuss both opportunities and challenges, encouraging students to consider perspectives of various stakeholders (e.g., small businesses, multinational corporations).
}

% Slide: Income Inequality
\begin{frame}{Income Inequality}
  \begin{itemize}
    \item \textbf{Trends:} Growing disparities in income and wealth globally and in Canada.
    \item \textbf{Causes:}
      \begin{itemize}
        \item Technological changes favoring skilled labor.
        \item Globalization and outsourcing.
        \item Declining unionization and policy shifts.
      \end{itemize}
    \item \textbf{Policy Responses:}
      \begin{itemize}
        \item Progressive taxation and social welfare programs.
        \item Education and workforce development.
      \end{itemize}
  \end{itemize}
\end{frame}

% Instructor Notes: Income Inequality
\note{
Discuss the Gini coefficient as a measure of inequality. Use Canadian data to illustrate trends. Pose the question: "Is rising inequality inevitable, or can policy effectively address it?"
}

% Slide: Climate Change Economics
\begin{frame}{Climate Change Economics}
  \begin{itemize}
    \item \textbf{Global Challenge:} Reducing emissions while supporting economic growth.
    \item \textbf{Canada’s Role:}
      \begin{itemize}
        \item Transition to renewable energy.
        \item Policies: Carbon pricing, green technology investments.
      \end{itemize}
    \item \textbf{Economic Opportunities:}
      \begin{itemize}
        \item Job creation in renewable sectors.
        \item Leading global innovation in sustainable technologies.
      \end{itemize}
  \end{itemize}
\end{frame}

% Instructor Notes: Climate Change Economics
\note{
Highlight carbon pricing as a Canadian success story. Discuss controversies around pipelines and fossil fuels, balancing economic and environmental goals. Encourage debate: "Should economic growth always take precedence over environmental concerns?"
}

% Slide: Technological Disruption
\begin{frame}{Technological Disruption}
  \begin{itemize}
    \item \textbf{Automation and AI:}
      \begin{itemize}
        \item Potential for increased productivity.
        \item Risks of job displacement in traditional industries.
      \end{itemize}
    \item \textbf{Implications for Policy:}
      \begin{itemize}
        \item Need for reskilling and education programs.
        \item Balancing innovation with ethical considerations.
      \end{itemize}
    \item \textbf{Global Competition:} Role of tech giants and innovation hubs.
  \end{itemize}
\end{frame}

% Instructor Notes: Technological Disruption
\note{
Discuss Canadian industries particularly affected by automation (e.g., manufacturing). Pose the question: "What skills will be most valuable in an AI-driven economy?"
}

% Slide: Post-Pandemic Recovery
\begin{frame}{Post-Pandemic Economic Recovery}
  \begin{itemize}
    \item \textbf{Challenges:}
      \begin{itemize}
        \item High inflation and debt levels.
        \item Supply chain disruptions.
      \end{itemize}
    \item \textbf{Opportunities:}
      \begin{itemize}
        \item Investment in healthcare and technology.
        \item Rethinking global supply chains.
      \end{itemize}
    \item \textbf{Policy Responses:}
      \begin{itemize}
        \item Monetary policy: Interest rates, quantitative easing.
        \item Fiscal policy: Stimulus packages and infrastructure investments.
      \end{itemize}
  \end{itemize}
\end{frame}

% Instructor Notes: Post-Pandemic Recovery
\note{
Discuss Canada’s fiscal stimulus and debt levels post-COVID. Compare with global approaches (e.g., U.S., EU). Pose the question: "How should Canada prioritize spending in the recovery phase?"
}

% Slide: Discussion Questions
\begin{frame}{Discussion Questions}
  \begin{itemize}
    \item How has globalization impacted Canada’s economy? Provide examples.
    \item What are the primary drivers of income inequality in Canada and globally?
    \item How can Canada balance economic growth with environmental sustainability?
  \end{itemize}
\end{frame}

% Instructor Notes: Discussion Questions
\note{
Encourage students to use specific examples to support their answers. For instance, globalization’s impact on Canadian manufacturing or climate policies like carbon taxes. Foster a debate on differing viewpoints.
}

% Slide: Critical Thinking Questions
\begin{frame}{Critical Thinking Questions}
  \begin{itemize}
    \item Can Canada remain competitive in a rapidly changing global economy? Why or why not?
    \item To what extent should Canada rely on natural resources for economic growth?
    \item What role does technology play in shaping future economic policies?
  \end{itemize}
\end{frame}

% Slide: Critical Thinking Questions and Answers
\begin{frame}{Critical Thinking Questions and Answers}
  \textbf{Question 1: Can Canada remain competitive in a rapidly changing global economy? Why or why not?}

  \textbf{Answer:}
  \begin{itemize}
    \item \textbf{Strengths:}
      \begin{itemize}
        \item Abundance of natural resources (e.g., oil, natural gas, minerals).
        \item Skilled and well-educated workforce.
        \item Stable political and financial systems attracting foreign investment.
      \end{itemize}
    \item \textbf{Challenges:}
      \begin{itemize}
        \item Global competition from emerging economies like China and India.
        \item Rapid technological disruption requiring adaptation.
        \item Need for environmental sustainability and reduced reliance on fossil fuels.
      \end{itemize}
    
  \end{itemize}
\end{frame}

% Slide: Critical Thinking Questions and Answers
\begin{frame}{Critical Thinking Questions and Answers}
  \textbf{Question 1: Can Canada remain competitive in a rapidly changing global economy? Why or why not?}

  \textbf{Answer:}
  \begin{itemize}

    \item \textbf{Policy Recommendations:}
      \begin{itemize}
        \item Invest in renewable energy and green technologies.
        \item Encourage innovation and R$\&$D in AI, biotechnology, and advanced manufacturing.
        \item Strengthen trade relationships and diversify exports.
      \end{itemize}
  \end{itemize}
\end{frame}






\begin{frame}{Critical Thinking Questions and Answers (cont'd)}
  \textbf{Question 2: To what extent should Canada rely on natural resources for economic growth?}
  \pause
  \textbf{Answer:}
  \begin{itemize}
    \item \textbf{Short-term reliance:}
      \begin{itemize}
        \item Natural resources drive economic stability, providing export revenues and jobs.
      \end{itemize}
    \item \textbf{Long-term risks:}
      \begin{itemize}
        \item Vulnerability to global commodity price fluctuations.
        \item Environmental concerns and global climate policy pressures.
      \end{itemize}
    \item \textbf{Path Forward:}
      \begin{itemize}
        \item Transition from resource-based growth to knowledge-based industries like technology and services.
        \item Invest in renewable energy and clean technologies.
        \item Support resource-dependent regions with diversification and reskilling policies.
      \end{itemize}
  \end{itemize}
\end{frame}

\begin{frame}{Critical Thinking Questions and Answers (cont'd)}
  \textbf{Question 3: What role does technology play in shaping future economic policies?}
  \pause
  \textbf{Answer:}
  \begin{itemize}
    \item \textbf{Opportunities:}
      \begin{itemize}
        \item Increases productivity through automation and AI.
        \item Creates new industries (e.g., green technology, fintech, and biotech).
        \item Enhances global connectivity for Canadian businesses.
      \end{itemize}
    \item \textbf{Challenges:}
      \begin{itemize}
        \item Job displacement in low-skilled industries.
        \item Regulation of technology, addressing concerns like privacy and cybersecurity.
        \item Bridging the digital divide for equitable access to technology.
      \end{itemize}
   
  \end{itemize}
\end{frame}


\begin{frame}{Critical Thinking Questions and Answers (cont'd)}
  \textbf{Question 3: What role does technology play in shaping future economic policies?}
  \pause
  \textbf{Answer:}
  \begin{itemize}

    \item \textbf{Policy Recommendations:}
      \begin{itemize}
        \item Invest in digital infrastructure and rural connectivity.
        \item Implement workforce reskilling programs to prepare workers for tech-driven industries.
        \item Collaborate with tech companies to regulate innovation ethically and sustainably.
      \end{itemize}
  \end{itemize}
\end{frame}

\end{document}
